
\documentclass{../notatki}

\title{Narzędzia matematyki w informatyce}

\begin{document}

Myślałbyś, drogi użytkowniku, że to podobny kaliber i jakość zajęć co mfi, lecz
nie! To ma sensowny wykład i ćwiczenia - triumf akademii.

\section{Teoria złożoności}

\begin{itemize}
  \item P - klasa problemów, które można rozwiązać w czasie wielomianowym $$
    c \cdot (|V| + |E|)^k
    $$
  \item NP - klasa problemów, których nie można rozwiązać w czasie wielomianowym
  \item NP-trudne - problemy, które można wielomianowo zamienić w
    inny problem w klasie NP.
  \item co-NP - klasa problemów, których komplement jest w klasie NP
\end{itemize}

Klasa P jest symetryczna. Jeśli problem decyzyjny należy do klasy P, to jego
komplement - negacja, również należy do klasy P. Klasa NP nie jest symetryczna.

\section{Grafy}

\begin{itemize}
  \item Graf regularny - wszystkie wierzchołki są tego samego stopnia
  \item Indukowany - podgraf stworzony przez wybranie podzbioru wierzchołków
  \item Składowe - spójne podgrafy, które nie są połączone
  \item Rozpinający - podgraf, który zawiera wszystkie wierzchołki
    grafu. \textbf{Niekoniecznie spójny} - kwestia konwencji
  \item Cykl Hamiltona - cykl, który przechodzi przez wszystkie wierzchołki
    grafu
\end{itemize}

Równość w grafie jest \textit{dziwna}. Aby grafy były równe, muszą mieć
identyczne zbiory wierzchołków i krawędzi. Jako że, elementy zbioru są
unikalne i rozróżnialne, to wpływ etykietowania na równość jest istotny.
Identyczne grafy, o innym etykietowaniu, są różne; można stworzyć nowy graf,
zmieniając etykietowanie.

\subsection{Graf Dwudzielny}

Graf dwudzielny, to graf, w którym wierzchołki są podzielone na dwa zbiory.
Między wierzchołkami z różnych zbiorów występują krawędzie, lecz
między wierzchołkami z tego samego zbioru nie ma krawędzi.

\begin{figure}[h]
  \centering
  \begin{tikzpicture}
    \node[draw, circle] (A) at (0,2) {A};
    \node[draw, circle] (B) at (0,1) {B};
    \node[draw, circle] (C) at (0,0) {C};
    \node[draw, circle] (D) at (4,2) {D};
    \node[draw, circle] (E) at (4,0) {E};
    \draw (A) -- (D);
    \draw (B) -- (D);
    \draw (C) -- (E);

    \node[
      draw,
      ellipse,
      fit=(A) (B) (C),
      inner sep=5pt,
    ] {};

    \node[
      draw,
      ellipse,
      fit=(D) (E),
      inner sep=5pt
    ] {};

  \end{tikzpicture}
  \caption{Przykład grafu dwudzielnego}
\end{figure}

Jeśli nie ma cyklu nieparzystego, i tylko wtedy, to graf jest dwudzielny.

\subsection{Drzewo}

Drzewo to graf acykliczny, w którym każdy wierzchołek ma jednego lub więcej
wierzchołków potomnych. Graf pusty, jest drzewem. Każde drzewo, jest lasem,
czyli niespójnym drzewem.

\subsection{Izomorfizm}

Izomorfizm to relacja między dwoma grafami, oznaczana $\cong$. Znaczy ona, że
istnieje bijekcja między wierzchołkami i krawędziami dwóch grafów. To
oznacza, że
aby grafy były izomorficzne, każdy wierzchołek i krawędź jednego grafu musi być
mapowany na odpowiedni wierzchołek i krawędź drugiego grafu,
zachowując relatywną strukturę.

To oznacza, że stopnie nie mogą się różnić, cykle muszą być identyczne, stopnie
sąsiedztwa muszą być identyczne, etc\dots \ . Nie ma prostego sposobu
na sprawdzenie.

\subsection{Skojarzenie}

Skojarzenie to zbiór rozłączny krawędzi, czyli podgraf, zdefiniowany
przez krawędzie, tak, że każdy wierzchołek jest końcem jednej
krawędzi. $\mu(G)$ to maksymalne skojarzenie w grafie $G$.
Skojarzenie doskonałe to takie, które jest maksymalnym skojarzeniem w
grafie $G$.

Mówimy, że wierzchołki należące do skojarzenia są nasycone.

\subsubsection{Twierdzenie Halla}

Graf dwudzielny $G = (V \cup W, E)$ ma skojarzenie sycące każdy wierzchołek
zbioru $W$ wtedy i tylko wtedy, gdy spełnione jest:
$$
\forall_{S \subseteq W} |N(S)| \geq |S|
$$
gdzie $N(S)$ to zbiór sąsiadów wierzchołków w $S$.

Tutaj istotnym jest obserwacja, że tym skojarzeniem, są połączenia między
dwoma zbiorami $V$ i $W$.

\subsubsection{Twierdzenie o małżeństwach}

Każdy $k$-regularny graf dwudzielny $G = (V \cup W, E)$ ma
skojarzenie doskonałe.

Kluczową obserwacją jest $|E| = k |V| = k |W|$. Z tego wynika, przy pomocy
twierdzenia Halla:
$$
\forall_{S \subseteq W} k|N(S)| \geq k|S| \quad \Rightarrow |S| \leq |N(S)|
$$

\begin{figure}[h]
  \centering
  \begin{tikzpicture}
    \node[draw, circle] (A) at (0,2) {};
    \node[draw, circle] (B) at (0,1) {};
    \node[draw, circle] (C) at (0,0) {};
    \node[draw, circle] (D) at (0,-1) {};
    \node[draw, circle] (E) at (0,-2) {};
    \node[draw, circle] (F) at (0,-3) {};

    \node[draw, circle] (G) at (4,2) {};
    \node[draw, circle] (H) at (4,1) {};
    \node[draw, circle] (I) at (4,0) {};
    \node[draw, circle] (J) at (4,-1) {};
    \node[draw, circle] (K) at (4,-2) {};
    \node[draw, circle] (L) at (4,-3) {};

    \draw (A) -- (H);
    \draw (A) -- (L);
    \draw (B) -- (J);
    \draw (B) -- (K);
    \draw (C) -- (I);
    \draw (C) -- (J);
    \draw[thick, color=red] (D) -- (G);
    \draw[thick, color=red] (D) -- (H);
    \draw[thick, color=red] (E) -- (H);
    \draw[thick, color=red] (E) -- (I);
    \draw (F) -- (L);
    \draw (F) -- (G);

    \node[
      draw,
      ellipse,
      fit=(A) (B) (C) (D) (E) (F),
      inner sep=5pt,
      label={[label distance=5pt]below:$V$}
    ] {};

    \node[
      draw,
      ellipse,
      fit=(G) (H) (I) (J) (K) (L),
      inner sep=5pt,
      label={[label distance=5pt]below:$W$}
    ] {};

    \node[
      draw,
      rectangle,
      fit=(D) (E),
      inner sep=3pt,
      label={[label distance=5pt]left:$S$}
    ] {};

    \node[
      draw,
      rectangle,
      fit=(G) (H) (I),
      inner sep=3pt,
      label={[label distance=5pt]right:$N(S)$}
    ] {};

  \end{tikzpicture}
  \caption{Graf z twierdzenia dla $k = 2$}
\end{figure}

\subsubsection{Twierdzenie Tutte}

Niech $G = (V, E)$ będzie dowolnym grafem. $G$ ma skojarzenie doskonałe wtedy i
tylko wtedy, gdy
$$
\forall_{S \subseteq V} o(G \setminus S) \leq |S|
$$
gdzie $o(G)$ to liczbą składowych nieparzystych w grafie $G$.

\end{document}
